\scanpage{117}
% line head | text | 66
\mspage{66}
% line 1–3 | text | Nous serons plus intéressés
\unsure{Nous} serons plus intéressés par la \deleted{\unsure{problème}}\inserted{situation} inverse,
\deleted{\unsure{d'écrire}} par la question de déterminer les
foncteurs
% line 4 | display | X, X subscript U,T,H, X prime
\[
\underset{\substack{\shortparallel\\X_{U,T,H}}}{X}\longrightarrow X'
\]
% line 5–7 | text | de X, obtenu par
de $X$, obtenu par $(U,T,H)$, dans un objet
quelconque de $\mathrm{Cat}$. On trouve tautologiquement

% line 8–10 | text | Proposition 3
\uline{Proposition 3}\quad $\operatorname{Hom}_{\mathrm{Cat}}(X,X')$ est en corr. bijective
canonique, fonctorielle en $X'$ dans $\mathrm{Cat}$, avec
l'ens des triples
% line 11–14 | display | f_U, f_T, u and the functor category
\[
\left\{\begin{aligned}
&\left.\begin{aligned}
f_U &: U\to X'\\
f_T &: T\to X'
\end{aligned}\right\}\ \text{flèches \unsure{dans|de} $\mathrm{Cat}$}\\
&u:H\longrightarrow H_{X'}(f_U^{\mathrm o}\times f_T)\quad\text{flèche \unsure{dans|de}}\\
&\hspace{9em}\underline{\operatorname{Hom}}(U^{\mathrm o}\times T,\mathrm{Ens})
\end{aligned}\right.
\]
% line 15 | text | où pour tt X prime
où pour \unsure{tt} $X'$ dans $\mathrm{Cat}$, on a posé
% line 16 | display | H of X prime and Hom
\[
H_{X'}=\bigl((x,y)\unsure{\mapsto|\vdash}\operatorname{Hom}_{X'}(x,y)\bigr):
{X'}^{\mathrm o}\times X'\longrightarrow\mathrm{Ens}.
\]

% line 17–18 | text | Sûrement on aura
\unsure{Sûrement} on \unsure{aura} $=$ un \deleted{équivalence}\inserted{isomorphisme} de
catégories
% line 19–20 | display | Underlined Hom and category of triples
\[
\begin{aligned}
\underline{\operatorname{Hom}}(X,X')&\simeq\text{catégorie des triples }(f_U,f_T,u)\\[-2pt]
&\hspace{3em}\text{comme \unsure{dans $(*)$|dessus}}.
\end{aligned}
\]

% line 21–22 | text | 1.4. Construction de sommes amalgamées
\uline{1.4}.\quad \uline{Construction de sommes amalgamées} \unsure{à} partir
\uline{d'une immersion ouverte ou fermée}.

% line 23–24 | text | Pour une construction
Pour une construction d'une \inserted{sorte de} somme
amalgamée, en termes d'un triple $(X,U,T)+\ldots$,
% archive sheet number: 117
